\documentclass[11pt,a4paper]{article}

\usepackage{fontspec}
\setmainfont{DejaVu Serif}[
  BoldFont={DejaVu Serif Bold},
  ItalicFont={DejaVu Serif},
  ItalicFeatures={FakeSlant=0.18},
  BoldItalicFont={DejaVu Serif Bold},
  BoldItalicFeatures={FakeSlant=0.18}
]
\setsansfont{DejaVu Sans}
\usepackage{polyglossia}
\setmainlanguage{english}
\usepackage[a4paper,margin=20.5mm]{geometry}
\usepackage{microtype}
\usepackage{setspace}
\setstretch{1.035}
\usepackage{amsmath,amssymb,amsthm,mathtools,bm}
\usepackage{booktabs,longtable,array,tabularx}
\usepackage{enumitem}
\usepackage{xcolor}
\usepackage{tikz}
\usetikzlibrary{arrows.meta,positioning,calc}
\usepackage[unicode,hidelinks]{hyperref}
\usepackage[nameinlink,noabbrev]{cleveref}
\emergencystretch=3em
\allowdisplaybreaks

\hypersetup{
  pdftitle={Frames, Forces, Constraints, and Dissipation under Observation Maps},
  pdfauthor={Stas, Independent Research Program},
  pdfsubject={A typed mechanics interface for Geometry of Observation},
  pdfkeywords={reference frames, Lagrange-d'Alembert, constraints, contact, friction, dissipation, observation geometry}
}

\definecolor{obsblue}{RGB}{37,83,138}
\definecolor{obsred}{RGB}{150,45,45}
\definecolor{obsgreen}{RGB}{30,115,78}
\definecolor{lightblue}{RGB}{239,246,253}
\definecolor{lightred}{RGB}{253,242,242}
\definecolor{lightgreen}{RGB}{239,249,244}

\theoremstyle{definition}
\newtheorem{definition}{Definition}[section]
\newtheorem{model}[definition]{Model}
\theoremstyle{plain}
\newtheorem{theorem}[definition]{Theorem}
\newtheorem{proposition}[definition]{Proposition}
\newtheorem{corollary}[definition]{Corollary}
\theoremstyle{remark}
\newtheorem{remark}[definition]{Remark}
\newtheorem{warning}[definition]{Boundary}

\newcommand{\R}{\mathbb{R}}
\newcommand{\dd}{\mathop{}\!d}
\newcommand{\norm}[1]{\left\lVert #1\right\rVert}
\newcommand{\pair}[2]{\left\langle #1,#2\right\rangle}
\newcommand{\SO}{\mathrm{SO}}
\newcommand{\SE}{\mathrm{SE}}
\newcommand{\Ocal}{\mathcal O}
\newcommand{\Kcal}{\mathcal K}
\newcommand{\Rcal}{\mathcal R}
\newcommand{\Qcal}{\mathcal Q}
\newcommand{\Acal}{\mathcal A}
\newcommand{\Tcal}{\mathcal T}
\newcommand{\Dcal}{\mathcal D}
\newcommand{\crossop}[1]{\left[#1\right]_\times}
\newcolumntype{L}[1]{>{\raggedright\arraybackslash}p{#1}}
\newcolumntype{Y}{>{\raggedright\arraybackslash}X}

\title{\bfseries Frames, Forces, Constraints, and Dissipation\\
\large A typed mechanics interface under observation maps}
\author{Stas\\
\small Independent Research Program; ORCID: 0009-0000-2294-705X}
\date{GO Core mechanics interface v0.1 --- 28 July 2026}

\begin{document}
\maketitle

\begin{center}
\fcolorbox{obsblue}{lightblue}{%
\parbox{0.92\linewidth}{\small
\textbf{Status.} Normative infrastructure for GO Core v0.2. This
interface fixes frame direction, point--vector separation, constraint
signs, force/power duality, and the nonnegative convention for
dissipation. It supplies a common contract for the Foucault,
bobsleigh, and roller-coaster modules. It does not propose a new law
of mechanics.}}
\end{center}

\begin{abstract}
Mechanical observation models become ambiguous when positions,
displacements, forces, generalized forces, and accelerations are added
without a common frame or when ``dissipated power'' is assigned both
signs in the same document. We define a typed interface on a
configuration manifold \(Q\). A frame edge
\(A_{F\leftarrow G}(t)\in\SO(3)\) maps free-vector components from
\(G\) to \(F\), while a point map additionally carries an origin.
Differentiating the affine point map yields the translational,
Coriolis, Euler, and centrifugal terms; these are representation
terms, not new interactions. Physical forces are covectors through
virtual power, and generalized forces belong to \(T_q^*Q\). Equality
constraints, unilateral contact, and Coulomb-type tangential reactions
are given separate signatures. Dissipation is defined by
\[
P_{\rm diss}:=-\pair{Q_d}{\dot q}\ge0,
\]
so a quadratic Rayleigh function produces
\(P_{\rm diss}=2\Rcal\). The Lagrangian energy identity is
\[
\dot E_L=P_{\rm act}-P_{\rm diss}+P_{\rm con}-\partial_tL.
\]
For an ideal scleronomous constraint \(P_{\rm con}=0\); for a moving
constraint it need not vanish. The interface also separates physical
frame changes, deterministic observation maps, stochastic
measurement channels, and estimators. Its statements are standard
mechanics recast as a machine-checkable contract.
\end{abstract}

\tableofcontents

\section{Scope and typed objects}

Let \(Q\) be a finite-dimensional configuration manifold and let
\(\mathsf F\) be a set of declared Euclidean frames. The interface
distinguishes the following objects.

\begin{center}
\begin{tabularx}{\linewidth}{L{31mm}L{35mm}Y}
\toprule
Object & Mathematical type & Admissible operation \\
\midrule
Point \(p\) & affine-space element & affine frame map, point difference \\
Free vector \(v\) & \(V_F\simeq\R^3\) & linear frame map, norm, dot product \\
Configuration velocity & \(\dot q\in T_qQ\) & pairing with a covector \\
Generalized force & \(Q_f\in T_q^*Q\) & virtual power \(\pair{Q_f}{\dot q}\) \\
Spatial force & \(F^F\in V_F\) & dot product with a coframed velocity \\
Acceleration & \(a^F\in V_F\) & addition only in one declared frame \\
Observation & map or Markov kernel & may be non-injective and lossy \\
\bottomrule
\end{tabularx}
\end{center}

\begin{warning}[Type firewall]
A frame transform is an invertible change of representation. An
observation map may be many-to-one. Neither is a force law. A
generalized force is a covector and cannot be added directly to a
spatial acceleration. Multiplication by a mass metric or an explicit
pullback is required.
\end{warning}

\section{Frame graph and affine point maps}

\begin{definition}[Directed frame edge]
For frames \(F,G\), the matrix
\[
A_{F\leftarrow G}(t)\in\SO(3)
\]
maps components of a free vector from \(G\) to \(F\):
\(v^F=A_{F\leftarrow G}v^G\). It obeys
\[
A_{F\leftarrow G}^{-1}=A_{G\leftarrow F}
=A_{F\leftarrow G}^{\mathsf T},\qquad
A_{F\leftarrow H}=A_{F\leftarrow G}A_{G\leftarrow H}.
\]
\end{definition}

A point needs an origin in addition to a rotation. Fix an inertial
frame \(I\), let \(c^I_F(t)\) be the origin of \(F\) represented in
\(I\), and abbreviate \(A(t)=A_{I\leftarrow F}(t)\). Then
\begin{equation}\label{eq:affine-point}
x^I=c^I_F+A\,x^F.
\end{equation}
The translation is forbidden for a free vector:
\(v^I=A v^F\), not \(c^I_F+A v^F\).

\begin{definition}[Angular velocity convention]
The \(F\)-components of angular velocity are defined by
\[
A^{\mathsf T}\dot A=\crossop{\Omega^F},\qquad
\crossop{\Omega}u=\Omega\times u.
\]
Thus \(\dot A=A\crossop{\Omega^F}\).
\end{definition}

\begin{theorem}[Kinematics of a moving frame]\label{thm:frame-kinematics}
If \(x^F(t)\) is twice differentiable, then
\begin{align}
v^I&=\dot c_F^I+
A\bigl(v^F+\Omega^F\times x^F\bigr),\label{eq:velocity-map}\\
a^I&=\ddot c_F^I+
A\bigl(
a^F+2\Omega^F\times v^F+\dot\Omega^F\times x^F
+\Omega^F\times(\Omega^F\times x^F)
\bigr).\label{eq:acceleration-map}
\end{align}
\end{theorem}

\begin{proof}
Differentiate \cref{eq:affine-point}, substitute
\(\dot A=A\crossop{\Omega^F}\), and differentiate once more. The product
rule gives two equal cross terms
\(\Omega^F\times v^F\), producing the coefficient \(2\).
\end{proof}

For a particle of mass \(m\) with physical force
\(F^F=A^{\mathsf T}F^I\), Newton's equation in \(F\) is
\begin{align}
m a^F={}&F^F
-mA^{\mathsf T}\ddot c_F^I
-2m\Omega^F\times v^F \notag\\
&-m\dot\Omega^F\times x^F
-m\Omega^F\times(\Omega^F\times x^F).
\label{eq:rotating-newton}
\end{align}
The last four terms are respectively translational, Coriolis, Euler,
and centrifugal effective forces.

\begin{warning}[Interpretation]
The effective terms in \cref{eq:rotating-newton} are produced by the
time-dependent representation. They must not be inserted in the
inertial equation and then transformed again.
\end{warning}

\section{Configuration embeddings and constraints}

Let a mechanical realization be a possibly time-dependent embedding
\[
\Phi_t:Q\longrightarrow(\R^3)^n,\qquad
q\longmapsto(x_1(t,q),\ldots,x_n(t,q)),
\]
with masses \(m_a>0\). For a time-independent embedding, the pullback
mass form is
\begin{equation}\label{eq:mass-metric}
G_{ij}(q)=\sum_{a=1}^n
m_a\,\partial_i x_a\cdot\partial_j x_a.
\end{equation}

\begin{proposition}[Mass-form rank]
The matrix \(G(q)\) in \cref{eq:mass-metric} is positive
semidefinite. It is positive definite exactly when
\(\dd\Phi_q:T_qQ\to(\R^3)^n\) is injective.
\end{proposition}

\begin{proof}
For \(\xi\in T_qQ\),
\[
\xi^{\mathsf T}G\xi
=\sum_a m_a\norm{\dd\Phi_q(\xi)_a}^2\ge0.
\]
It vanishes exactly on \(\ker\dd\Phi_q\).
\end{proof}

Equality constraints are written
\[
h(t,q)=0,\qquad A_h(t,q)=\partial_qh.
\]
For a regular ideal constraint with multiplier \(\lambda\), the
reaction covector is
\[
Q_h=A_h^{\mathsf T}\lambda.
\]
An admissible velocity satisfies
\begin{equation}\label{eq:rheo-velocity}
A_h\dot q+\partial_t h=0.
\end{equation}
Therefore its reaction power is
\begin{equation}\label{eq:rheo-power}
P_h=\pair{Q_h}{\dot q}
=-\lambda^{\mathsf T}\partial_t h.
\end{equation}
It vanishes for a scleronomous constraint
\(\partial_t h=0\), but not in general for a moving support.

For unilateral contact we fix the sign convention
\begin{equation}\label{eq:complementarity}
g_j(t,q)\ge0,\qquad N_j\ge0,\qquad N_jg_j=0,
\end{equation}
where \(g_j>0\) means separation and \(N_j\) is compressive.
Impacts, restitution, and measure-valued impulses require a separate
nonsmooth extension and are outside the smooth interface.

\section{Forces, pullback, and virtual power}

Suppose the spatial velocity of application point \(a\) is
\[
v_a=J_a(t,q)\dot q+v_{a,0}(t,q).
\]
A spatial force \(F_a\) pulls back to
\[
Q_a=J_a^{\mathsf T}F_a\in T_q^*Q.
\]
For the virtual part,
\[
\pair{Q_a}{\delta q}=F_a\cdot J_a\delta q.
\]
The actual spatial power contains the affine contribution:
\[
F_a\cdot v_a
=\pair{Q_a}{\dot q}+F_a\cdot v_{a,0}.
\]
Consequently a time-dependent embedding cannot be treated as a
time-independent Jacobian without losing support power.

\begin{proposition}[Orthogonal covariance of power]
For \(A\in\SO(3)\), if \(F^I=AF^F\) and \(v^I=Av^F\), then
\[
F^I\cdot v^I=F^F\cdot v^F.
\]
\end{proposition}

\begin{proof}
\((AF^F)^{\mathsf T}(Av^F)
=(F^F)^{\mathsf T}A^{\mathsf T}Av^F\).
\end{proof}

This proposition applies to coframed physical vectors. It does not
authorize pairing a force in one frame with a velocity in another.

\section{Dissipation and contact friction}

\begin{definition}[Dissipative generalized force]
A generalized force \(Q_d(t,q,\dot q)\) is dissipative on an
admissible domain if
\[
\pair{Q_d}{\dot q}\le0.
\]
Its nonnegative dissipated power is
\begin{equation}\label{eq:pdiss}
P_{\rm diss}:=-\pair{Q_d}{\dot q}\ge0.
\end{equation}
\end{definition}

For a Rayleigh function
\[
\Rcal(t,q,\dot q)=\frac12\dot q^{\mathsf T}C(t,q)\dot q,
\qquad C=C^{\mathsf T}\succeq0,
\]
the convention is
\[
Q_d=-\partial_{\dot q}\Rcal=-C\dot q,\qquad
P_{\rm diss}=\dot q^{\mathsf T}C\dot q=2\Rcal.
\]

At contact \(j\), let \(n_j\) be the unit normal and
\(P_{T,j}=I-n_jn_j^{\mathsf T}\). Write the force on the modeled body
as
\[
F_j=N_jn_j+T_j,\qquad T_j\cdot n_j=0,
\]
and the tangential relative velocity as
\[
v_{T,j}=P_{T,j}(v_{\rm body,j}-v_{\rm support,j}).
\]
A Coulomb-type admissible reaction obeys
\begin{equation}\label{eq:coulomb}
\norm{T_j}\le\mu_jN_j,\qquad
T_j\cdot v_{T,j}\le0.
\end{equation}
Thus
\begin{equation}\label{eq:friction-power}
P_{{\rm fric},j}=-T_j\cdot v_{T,j}\ge0.
\end{equation}
For sliding with \(v_{T,j}\ne0\), the maximum-dissipation selection is
\[
T_j=-\mu_jN_j\,\frac{v_{T,j}}{\norm{v_{T,j}}}.
\]
For sticking, \(v_{T,j}=0\), the friction force is set-valued inside
the cone and its instantaneous power is zero.

\begin{definition}[Guarded cone utilization]
On an admissible contact state define
\[
u_j=
\begin{cases}
\norm{T_j}/(\mu_jN_j),&\mu_jN_j>0,\\
0,&\mu_jN_j=0\ \text{and}\ T_j=0.
\end{cases}
\]
The cone law excludes the missing case
\(\mu_jN_j=0,\ T_j\ne0\). For a finite contact set,
\[
S_{\rm slip}=1-\max_j u_j\in[0,1],
\]
with \(\max\varnothing:=0\).
\end{definition}

\section{Lagrange--d'Alembert energy identity}

Let \(L(t,q,\dot q)\) be regular on the active smooth regime and set
\[
E_L=\pair{\partial_{\dot q}L}{\dot q}-L.
\]
The forced equations are
\begin{equation}\label{eq:lda}
\frac{\dd}{\dd t}\partial_{\dot q}L-\partial_qL
=Q_{\rm act}+Q_d+Q_{\rm con}.
\end{equation}

\begin{theorem}[Energy balance]\label{thm:energy-balance}
Every \(C^2\) solution of \cref{eq:lda} satisfies
\begin{equation}\label{eq:energy-balance}
\dot E_L
=P_{\rm act}-P_{\rm diss}+P_{\rm con}-\partial_tL,
\end{equation}
where
\[
P_{\rm act}=\pair{Q_{\rm act}}{\dot q},\qquad
P_{\rm con}=\pair{Q_{\rm con}}{\dot q}.
\]
\end{theorem}

\begin{proof}
Direct differentiation gives
\[
\dot E_L
=\pair{\frac{\dd}{\dd t}\partial_{\dot q}L-\partial_qL}{\dot q}
-\partial_tL.
\]
Insert \cref{eq:lda} and use
\cref{eq:pdiss}.
\end{proof}

\begin{corollary}[Autonomous ideal case]
If \(\partial_tL=0\), \(Q_{\rm act}=0\),
\(Q_d=0\), and the constraints are ideal and scleronomous, then
\(\dot E_L=0\).
\end{corollary}

\begin{warning}[No-double-counting rule]
A moving track may be represented either by an ambient formulation
with explicit reaction/support power or by reduction to a
time-dependent Lagrangian. The same support work must not be added in
both places.
\end{warning}

\section{Observation protocol and transformation classes}

A complete protocol is the typed chain
\[
X_t\xrightarrow{\Rcal_{\rm red}}\widetilde X_t
\xrightarrow{\Ocal_\lambda}Z_t
\xrightarrow{\Kcal_\lambda}\operatorname{Law}(Y_t\mid Z_t)
\xrightarrow{\widehat\theta_\lambda}\widehat\theta_t.
\]
The protocol record \(\lambda\) includes:
\begin{enumerate}[label=(\roman*)]
  \item frame graph, clock, coordinates, and unit context;
  \item system descriptor and constraint convention;
  \item spatial and temporal resolution and observation horizon;
  \item deterministic readout and stochastic/noiseless declaration;
  \item calibration, estimator, uncertainty, and missing-data policy;
  \item a declared transformation group for every invariant claim.
\end{enumerate}

\begin{warning}[Invariant scope]
Rigid-motion invariance, covariance under a time-dependent frame,
Galilean invariance, label-swap invariance, and invariance under a
lossy observation map are different statements. The word
``invariant'' is rejected unless the acting group and the transformed
objects are named.
\end{warning}

\section{Machine gates}

The mechanics extension admits a reference document only if the
following gates are satisfied.

\begin{longtable}{L{31mm}L{46mm}L{65mm}}
\toprule
Gate & Required declaration & Rejected pattern \\
\midrule
\endhead
Frame closure & every vector sum has one frame; every point map has an
origin & \(r^I+q^B\) without \(A_{I\leftarrow B}\) \\
Force type & spatial force or generalized covector, with pullback map
& adding \(F\) directly to \(a\) \\
Constraint sign & gap side, reaction sign, complementarity, active set
& unsigned \(g,N\) convention \\
Dissipation sign & \(P_{\rm diss}\ge0\) and force-on-system convention
& \(\dot E=-\sum T\cdot v_{\rm rel}\) with \(T\cdot v_{\rm rel}\le0\) \\
Moving support & explicit \(\partial_t h\), support power, or reduced
time-dependent \(L\) & declaring every ideal constraint workless \\
Friction guard & zero-load branch for \(\norm T/(\mu N)\) & division by
\(\mu N=0\) \\
Observation split & frame map, observation, channel, estimator kept
separate & calling a camera projection a frame transform \\
Invariant scope & named group and hypotheses & unqualified
``observation invariant'' \\
\bottomrule
\end{longtable}

\section{Claim boundary}

The interface is a consistency layer. It does not establish:
\begin{itemize}
  \item existence or uniqueness for arbitrary nonsmooth contact;
  \item a constitutive law for ice, wheel--rail, or air resistance;
  \item structural safety, fatigue, or human-tolerance limits;
  \item a precision geophysical pendulum model;
  \item scientific novelty of the underlying standard identities.
\end{itemize}
Any application must add its own regularity, constitutive, parameter,
and uncertainty hypotheses.

\begin{thebibliography}{9}

\bibitem{bloch}
A.~M. Bloch, \emph{Nonholonomic Mechanics and Control}, 2nd ed.,
Springer, 2015. \href{https://doi.org/10.1007/978-1-4939-3017-3}
{doi:10.1007/978-1-4939-3017-3}.

\bibitem{brogliato}
B. Brogliato, \emph{Nonsmooth Mechanics: Models, Dynamics and
Control}, 3rd ed., Springer, 2016.
\href{https://doi.org/10.1007/978-3-319-28664-8}
{doi:10.1007/978-3-319-28664-8}.

\bibitem{marsden}
J.~E. Marsden and T.~S. Ratiu,
\emph{Introduction to Mechanics and Symmetry}, 2nd ed., Springer,
1999. \href{https://doi.org/10.1007/978-0-387-21792-5}
{doi:10.1007/978-0-387-21792-5}.

\bibitem{arnold}
V.~I. Arnold, \emph{Mathematical Methods of Classical Mechanics},
2nd ed., Springer, 1989.
\href{https://doi.org/10.1007/978-1-4757-2063-1}
{doi:10.1007/978-1-4757-2063-1}.

\end{thebibliography}

\end{document}
